\documentclass{article}
\usepackage{style}

\title{Experiment 4 - Cosmic ray detection}
\author{Daire O'Brien - 24365676}
\begin{document}
\maketitle
\section{Background}
    In 1910 German physicist Theodore Wulf, using a early form of ionized gas detector called an electrometer, discovered that the amount of radiation at the top of the Eiffel Tower was noticibly higher then at the bottom (Th. Wulf, 1910)$^{[1]}$. Using balloons to achive altitudes of around $\qty{9}{\km}$, other physicists found the rate of ionization was higher at higher altitudes. This led physicists to conclude that there was some source of radiation not coming from the Earth. \\
    Later work by Bruno Rossi found that most of the rays consisted of alpha particles. Additionally, he also found that two of his Geiger counters triggered in a time span shorter then expected, especially across longer distances. This led him to conclude that the events he was detecting were broader showers of particles. (Bruno Rossi, 1934)$^{[2]}$.
    Cosmic rays are high energy particles emitted from high energy events from local, galatic or intergalactic space. They typically consist of protons or small atomic nuclei. 
\section{References}
\begin{enumerate}
    \item Wulf, Theodore, 1910, \emph{Beobachtungen \"uber die Strahlung hoher Durchdringgungsf\"ahigkeit auf dem Eiffelturm}, Physikalische Zeitschrift 11.1910. Accessed through https://digitalisate.sub.uni-hamburg.de/recherche
    \item Rossi, Bruno, 1934, \emph{Misure sulla distribuzione angolare di intensita della radiazione penetrante all'Asmara}, Ricerca Scientifica, vol. 5
\end{enumerate}
\end{document}